\documentclass[spanish,letterpaper,oneside]{article}
\usepackage[utf8]{inputenc}
\usepackage[spanish]{babel}
\usepackage{geometry}
\geometry{letterpaper, margin=2.5cm}
\usepackage{xcolor}
\usepackage{enumitem}
\usepackage{tcolorbox}

\begin{document}
\begin{titlepage}
    \centering
    \vspace*{2cm}
    {\Huge \textbf{Universidad Autónoma de Coahuila} \par}
    \vspace{1cm}
    {\Large Facultad de Medicina \par}
    \vspace{2cm}
    {\LARGE \textbf{Examen de Ingreso: Salud Pública y Medicina Comunitaria (Premedicina)} \par}
    \vspace{1cm}
    {\Large Guía Autocontenida y Banco de Reactivos EXANI-II \par}
    \vspace{3cm}
    {\large Elaborado por: Aspirante a Médico Cirujano \par}
    \vfill
    {\large \today \par}
\end{titlepage}

\tableofcontents
\newpage

\section{Módulo Teórico: Epidemiología y Prevención (Alto Rendimiento)}
\subsection{El proceso de Salud-Enfermedad (Tríada Ecológica)}
\begin{itemize}
    \item \textbf{Huésped:} Cualquier organismo vivo capaz de ser infectado (albergar un agente). Depende de la edad, sexo, estado inmune, estado nutricional.
    \item \textbf{Agente:} Todo ente u organismo cuya presencia (o ausencia) es imprescindible para desencadenar la enfermedad. (Biológicos [virus/bacterias], Físicos [radiación, calor], Químicos [venenos]).
    \item \textbf{Ambiente:} Factores físicos (clima, geografía), biológicos (flora, fauna, vectores) o socioculturales que facilitan el encuentro entre agente y huésped. 
    \item \textbf{Vector:} Ser vivo (ej. mosquito, garrapata) portador del agente y que lo transmite al huésped.
    \item \textbf{Fómite:} Objeto inanimado (ej. jeringa, picaporte) contaminado capaz de transmitir el agente.
\end{itemize}

\subsection{Historia Natural de la Enfermedad (Modelo de Leavell y Clark)}
Este paradigma divide el curso de toda enfermedad en ausencia de intervención médica:
\begin{enumerate}
    \item \textbf{Período Prepatogénico:} Ocurre antes de la enfermedad. El huésped, agente y ambiente interactúan. Si se rompe el equilibrio, ocurre el contagio. Aquí actúa la \textit{Prevención Primaria}.
    \item \textbf{Período Patogénico:} Ya existe la enfermedad (cambios celulares/tisulares).
    \begin{itemize}
        \item \textit{Fase subclínica (incubación):} Cambios silenciosos. Aún no hay síntomas.
        \item \textit{Horizonte clínico:} Cruce de la línea de visibilidad donde aparecen síntomas. Aquí actúa la \textit{Prevención Secundaria}.
        \item \textit{Resolución:} Muerte, cronicidad, incapacidad o curación. (Incapacidad crónica requiere \textit{Prevención Terciaria}).
    \end{itemize}
\end{enumerate}

\subsection{Niveles de Prevención (EXANI-II)}
\begin{tcolorbox}[colback=orange!5!white,colframe=orange!80!black,title=Los 3 Niveles Prevencionistas Clásicos]
\textbf{1. Prevención Primaria:} Intervenir \underline{antes} de enfermar (Prepatogénico). Dos estrategias: Fomento a la salud (educación sexual, nutrición) y Protección específica (Vacunación, uso de condón).\\
\textbf{2. Prevención Secundaria:} Intervenir de manera \underline{temprana}. Dos estrategias: Diagnóstico precoz y Tratamiento oportuno (Tamizajes, Papanicolau, Mastografías) para detener el progreso.\\
\textbf{3. Prevención Terciaria:} Rehabilitación (física, psicológica y social) frente a secuelas o discapacidad, buscando reintegrar al paciente a la sociedad (ej. prótesis, terapia física post-EVC).
\end{tcolorbox}

\subsection{Signos vs Síntomas / Epidemiología}
\begin{itemize}
    \item \textbf{Signo (Objetivo):} Medible, cuantificable, observable por el médico. Ej. Fiebre (39°C), taquicardia (120 lpm), eritema, edema.
    \item \textbf{Síntoma (Subjetivo):} Percepción experimentada por el paciente, no comprobable matemáticamente. Ej. Dolor (cefalea, mialgia), astenia (cansancio), adinamia (falta de fuerza), mareo.
    \item \textbf{Síndrome:} Conjunto organizado de signos y síntomas con causa variable y origen patofisiológico común (ej. Síndrome Febril, Síndrome Metabólico).
    \item \textbf{Incidencia:} Número de casos NUEVOS de una enfermedad en un periodo.
    \item \textbf{Prevalencia:} Número TOTAL (nuevos + acumulados viejos) de casos.
    \item \textbf{Endemia:} Enfermedad habitual constante en una zona geográfica (ej. Malaria en África).
    \item \textbf{Epidemia / Brote:} Aumento súbito de casos por encima de lo esperado. Pandemia si cruza fronteras continentales.
\end{itemize}

\newpage
\section{Banco de Reactivos: Simulación EXANI-II}
\begin{enumerate}[label=\textbf{\arabic*.}]
    \item \textbf{¿A qué componente de la tríada ecológica corresponde el estado de desnutrición de un infante que facilita el desarrollo de una enfermedad diarreica?}
    \begin{enumerate}
        \item[A)] Ambiente socioeconómico
        \item[B)] Factor predisponente del Agente
        \item[C)] Características del Huésped
        \item[D)] Inmunidad adaptativa
    \end{enumerate}
    \textit{Respuesta: C. La desnutrición es una característica inherente al organismo vivo (huésped) que altera su susceptibilidad a la infección.}

    \item \textbf{Un paciente contrae malaria (paludismo). El mosquito Anopheles hembra que le picó es considerado epidemiológicamente como:}
    \begin{enumerate}
        \item[A)] Fómite
        \item[B)] Vector biológico (Agente transmisor del ambiente)
        \item[C)] Agente causal directo
        \item[D)] Huésped definitivo
    \end{enumerate}
    \textit{Respuesta: B. El mosquito es el vector (vehículo vivo transmisor). El agente causal real es el parásito (Plasmodium sp.).}

    \item \textbf{La aplicación sistemática del Papanicolau (citología cervical) a mujeres sanas para detectar células precancerosas del cérvix, se clasifica como una acción de prevención:}
    \begin{enumerate}
        \item[A)] Primaria (Protección específica)
        \item[B)] Secundaria (Diagnóstico precoz)
        \item[C)] Terciaria (Rehabilitación)
        \item[D)] Cuaternaria
    \end{enumerate}
    \textit{Respuesta: B. El Papanicolau no "evita" que el VPH contagie (eso lo hace la vacuna, Prev. 1a); el tamizaje es diagnóstico precoz (Prev. 2a) detectando la enfermedad cuando aún es subclínica y curable.}

    \item \textbf{Durante la exploración física, el médico documenta presión arterial de 150/95 mmHg y anota un soplo cardíaco. El paciente, por otro lado, refiere tener "dolor en el pecho y mucha fatiga al caminar". ¿Cuál de las opciones representa un SÍNTOMA?}
    \begin{enumerate}
        \item[A)] Presión arterial 150/95
        \item[B)] Soplo cardíaco
        \item[C)] Taquicardia
        \item[D)] Fatiga (astenia) y dolor
    \end{enumerate}
    \textit{Respuesta: D. La fatiga y el dolor son subjetivos (experiencia del paciente). La presión y los soplos son objetivos y medibles (Signos).}

    \item \textbf{En epidemiología, la tasa que expresa la frecuencia de casos NUEVOS de COVID-19 en una ciudad durante la última semana, excluyendo a los que ya estaban enfermos desde hace un mes, se denomina:}
    \begin{enumerate}
        \item[A)] Prevalencia
        \item[B)] Letalidad
        \item[C)] Mortalidad específica
        \item[D)] Incidencia
    \end{enumerate}
    \textit{Respuesta: D. Incidencia = Nuevos. Prevalencia = Total (Nuevos + Antiguos acumulados). Letalidad = Riesgo de morir si enfermas.}

    \item \textbf{El desarrollo de campañas publicitarias en televisión orientadas a recomendar la disminución del consumo de bebidas azucaradas, pertenece al nivel de Prevención Primaria, específicamente al rubro de:}
    \begin{enumerate}
        \item[A)] Protección Específica
        \item[B)] Promoción / Fomento a la Salud
        \item[C)] Tamizaje poblacional
        \item[D)] Rehabilitación nutricional
    \end{enumerate}
    \textit{Respuesta: B. La promoción fomenta cambios conductuales y estilos de vida generales. La "Protección específica" implica barreras físicas o químicas directas contra una sola enfermedad (ej. casco, condón, vacuna).}

    \item \textbf{De acuerdo al modelo de Historia Natural de la Enfermedad, ¿qué suceso anatómico/fisiológico marca el paso entre la fase subclínica (incubación) y la fase clínica evidente?}
    \begin{enumerate}
        \item[A)] El Horizonte Clínico (aparición del primer signo/síntoma)
        \item[B)] La convalecencia
        \item[C)] El inicio del periodo prepatogénico
        \item[D)] El contagio
    \end{enumerate}
    \textit{Respuesta: A. El horizonte clínico es el umbral imaginario cruzado cuando el paciente nota la sintomatología de la enfermedad. El contagio inicia la incubación celular subclínica.}

    \item \textbf{Según la OMS (1948), el concepto moderno y positivo de salud es definido como:}
    \begin{enumerate}
        \item[A)] La ausencia absoluta de dolor y discapacidad física.
        \item[B)] Un completo estado de bienestar físico, mental y social, y no solo la ausencia de enfermedad.
        \item[C)] El equilibrio puramente biológico interno (homeostasis).
        \item[D)] La capacidad inmunológica para resistir infecciones estacionales.
    \end{enumerate}
    \textit{Respuesta: B. Esta definición oficial integró por primera vez el aspecto mental y social al bienestar del paciente.}

    \item \textbf{Un niño contrae tétanos al pisar un clavo oxidado tirado en un campo baldío. En esta cadena epidemiológica, el clavo actúa como:}
    \begin{enumerate}
        \item[A)] Fómite (Reservorio inanimado o vehículo pasivo)
        \item[B)] Vector mecánico vivo
        \item[C)] Huésped secundario
        \item[D)] Agente etiológico
    \end{enumerate}
    \textit{Respuesta: A. Un fómite es un objeto inerte y sin vida capaz de transmitir microorganismos (ej. jeringas, cubiertos, ropa, el clavo). El agente causal directo de la enfermedad es la toxina del Clostridium tetani.}

    \item \textbf{Dentro del modelo clínico, ¿cómo se denomina a la evolución probable o curso esperado de una enfermedad y su desenlace en el paciente, en función de su tratamiento?}
    \begin{enumerate}
        \item[A)] Diagnóstico presuntivo
        \item[B)] Pronóstico
        \item[C)] Etiología
        \item[D)] Patogenia
    \end{enumerate}
    \textit{Respuesta: B. El pronóstico es la "predicción" médica de supervivencia, curación o secuelas (ej. "pronóstico reservado", "favorable"). La etiología es la causa originaria de la enfermedad.}

    \item \textbf{En el esquema tradicional de vacunación infantil en México (y gran parte del mundo), las vacunas se aplican idealmente de forma programada a niños sanos. Este acto médico es el pilar de:}
    \begin{enumerate}
        \item[A)] Prevención primaria (Protección específica)
        \item[B)] Prevención secundaria
        \item[C)] Prevención terciaria
        \item[D)] Medicina paliativa
    \end{enumerate}
    \textit{Respuesta: A. Las vacunas inmunizan ANTES de que el paciente interactúe con el patógeno (periodo prepatogénico), blindándolo específicamente contra él (protección específica).}

    \item \textbf{Una enfermedad es diagnosticada como "Endémica" (Endemia) en una región si cumple la siguiente característica:}
    \begin{enumerate}
        \item[A)] Se diseminó rápidamente a más de un continente.
        \item[B)] Presenta una tasa de prevalencia estable y continua esperada a lo largo del tiempo en dicha zona geográfica.
        \item[C)] Ocurre un aumento súbito e inesperado de mortalidad a cero.
        \item[D)] Fue completamente erradicada de la zona hace 5 años.
    \end{enumerate}
    \textit{Respuesta: B. Endemia significa "propio del lugar", donde existe un número esperado y sostenido de casos anuales (ej. Chagas o Dengue en algunas regiones de Latinoamérica). Un "pico" súbito fuera de ese margen es una Epidemia/Brote.}

    \item \textbf{La colocación de una prótesis de cadera y posteriores 6 meses de terapia de reacondicionamiento en una persona mayor que sufrió una caída, es un ejemplo práctico de prevención:}
    \begin{enumerate}
        \item[A)] Primaria
        \item[B)] Secundaria (Limitar el daño)
        \item[C)] Terciaria (Rehabilitación)
        \item[D)] Cuaternaria
    \end{enumerate}
    \textit{Respuesta: C. La enfermedad traumática ya dejó su secuela destructiva. El uso de prótesis y terapia busca reintegrarlo a su actividad motora independiente.}

    \item \textbf{En la historia de la salud pública, la epidemia global de Gripe Española de 1918 y la reciente emergencia de SARS-CoV-2 (2020) que afectaron de manera descontrolada a varios continentes a la vez, son catalogadas como:}
    \begin{enumerate}
        \item[A)] Hiperendemias
        \item[B)] Epidemias locales convergentes
        \item[C)] Pandemias
        \item[D)] Sindemias
    \end{enumerate}
    \textit{Respuesta: C. Pandemia es la afectación epidémica que cruza fronteras internacionales afectando una vasta área global.}

    \item \textbf{¿Cómo se llama la vía de transmisión en la cual una mujer embarazada, infectada por sífilis o VIH, transmite el microorganismo a su feto a través de la placenta durante la gestación?}
    \begin{enumerate}
        \item[A)] Transmisión horizontal indirecta
        \item[B)] Transmisión vertical
        \item[C)] Transmisión iatrogénica
        \item[D)] Infección nosocomial
    \end{enumerate}
    \textit{Respuesta: B. La transmisión de madre a hijo (transplacentaria, durante el parto o por leche materna) se conoce como vertical.}

    \item \textbf{Un individuo expuesto a la toxina botulínica en una lata de comida en conserva no presenta síntomas durante las primeras 36 horas. En el modelo de Leavell y Clark, el paciente se encuentra transitando por:}
    \begin{enumerate}
        \item[A)] Período Prepatogénico
        \item[B)] Etapa Clínica (horizonte clínico)
        \item[C)] Fase Subclínica de Incubación/Latencia (Periodo Patogénico)
        \item[D)] Cronicidad y Rehabilitación
    \end{enumerate}
    \textit{Respuesta: C. El agente ya está alterando las rutas bioquímicas del paciente internamente, por lo que es patogénico, pero los síntomas aún no superan el horizonte clínico (no hay parálisis aún).}

    \item \textbf{El asbesto, radiación UV intensa, tabaco o una mala dieta alta en grasas poliinsaturadas oxidadas, todos ellos actúan en la cadena de causalidad de cáncer primario y se clasifican ecológicamente como agentes:}
    \begin{enumerate}
        \item[A)] Biológicos y vectoriales
        \item[B)] Físicos o Químicos carcinógenos
        \item[C)] Ambientales psicológicos
        \item[D)] Intrínsecos del huésped hereditario
    \end{enumerate}
    \textit{Respuesta: B. Son agentes patógenos de naturaleza fisicoquímica que influyen crónicamente. El agente no necesita ser un microbio.}

    \item \textbf{Indicador en salud que mide la cantidad de fallecimientos totales respecto a toda la población en un determinado periodo y región (tasa cruda):}
    \begin{enumerate}
        \item[A)] Tasa de Morbilidad
        \item[B)] Tasa de Letalidad
        \item[C)] Tasa de Mortalidad
        \item[D)] Esperanza de Vida
    \end{enumerate}
    \textit{Respuesta: C. Mortalidad es muerte poblacional. Morbilidad son los que "enferman". Letalidad es cuántos "enfermos" terminaron muertos por dicha enfermedad específica.}

    \item \textbf{Término médico-epidemiológico utilizado para describir una infección adquirida intrahospitalariamente (ej. un paciente internado por un infarto que a los 5 días contrae neumonía en el hospital):}
    \begin{enumerate}
        \item[A)] Infección comunitaria adquirida
        \item[B)] Zoonosis
        \item[C)] Enfermedad iatrogénica
        \item[D)] Infección Nosocomial (Infección asociada a la atención a la salud - IAAS)
    \end{enumerate}
    \textit{Respuesta: D. Toda infección que incuba o aparece a partir de las 48 hrs del ingreso a la instalación médica se llama nosocomial o IAAS.}

    \item \textbf{El aislamiento y cuarentena obligatoria (como se realizó durante los brotes de ébola o cólera aguda severa) son estrategias epidemiológicas enfocadas fundamentalmente a romper el eslabón de:}
    \begin{enumerate}
        \item[A)] Susceptibilidad del huésped sano
        \item[B)] La vía de transmisión desde la puerta de salida hacia los sanos
        \item[C)] Resistencia del vector
        \item[D)] Disminución de secuelas locales a largo plazo
    \end{enumerate}
    \textit{Respuesta: B. El aislamiento encierra al individuo ya infectado y su puerta de salida para cortar cualquier posible vehículo de transmisión al resto del ambiente/comunidad.}

\end{enumerate}
\end{document}
